\documentclass[12pt,a4paper]{article}

\usepackage{amsmath}
\usepackage{fontspec}
\setmainfont{Times New Roman}
\setmonofont{Consolas}[Scale=0.88]
\usepackage[a4paper,top=2.2cm,bottom=2.2cm,left=2.5cm,right=2.5cm]{geometry}
\usepackage{microtype}
\usepackage{setspace}
\setstretch{1.1}
\usepackage{array,tabularx,booktabs}
\usepackage{enumitem}
\usepackage{float}
\usepackage{titlesec}
\usepackage{fancyhdr}
\usepackage{xcolor}
\usepackage[hidelinks]{hyperref}

\definecolor{accent}{HTML}{1F3A5F}
\hypersetup{colorlinks=true,linkcolor=accent,urlcolor=accent}

\titleformat{\section}{\large\bfseries\color{accent}}{\thesection.}{0.6em}{}
\titleformat{\subsection}{\normalsize\bfseries\color{accent}}{\thesubsection.}{0.6em}{}
\titlespacing*{\section}{0pt}{1.4em}{0.6em}
\titlespacing*{\subsection}{0pt}{1.1em}{0.4em}

\setlist{itemsep=0.25em,topsep=0.4em}
\newcolumntype{L}{>{\raggedright\arraybackslash}X}
\renewcommand{\arraystretch}{1.15}
\renewcommand{\tablename}{Bảng}
\setlength{\parindent}{0pt}
\setlength{\parskip}{0.5em}

\pagestyle{fancy}
\fancyhf{}
\fancyhead[L]{\small\color{gray}Báo cáo nghiên cứu PH-SHOWOA}
\fancyhead[R]{\small\color{gray}Lam Huynh Khang}
\fancyfoot[C]{\small\thepage}
\renewcommand{\headrulewidth}{0.4pt}

\begin{document}

\begin{center}
  {\LARGE\bfseries\color{accent} BÁO CÁO NGHIÊN CỨU}\\[0.5em]
  {\large Cải tiến thuật toán lai PH-SHOWOA trên GPU\\cho bài toán VRPSPDTW}\\[1em]
  {\normalsize Lam Huynh Khang}\\[0.2em]
  {\small Ngày 24 tháng 10 năm 2026}
\end{center}
\vspace{0.5em}
\hrule
\vspace{0.5em}

\section{Đề tài nghiên cứu và kết quả đạt được}

Đề tài tập trung cải tiến thuật toán lai PH-SHOWOA (Spotted Hyena Optimizer kết hợp Whale Optimization Algorithm) cho bài toán VRPSPDTW theo hai hướng chính: chuyển toàn bộ vòng tiến hóa lên GPU và thay thế bước khởi tạo SA bằng quy trình SA\_RCRS\_GRASP. Đến thời điểm báo cáo, thuật toán gốc đã được tái lập hoàn chỉnh và phiên bản cải tiến đã được cài đặt; quá trình thực nghiệm đang được tiến hành, do đó chưa có kết quả chính thức.

\paragraph{Bài toán.} Vehicle Routing Problem with Simultaneous Pickup and Delivery and Time Windows (VRPSPDTW): mỗi khách hàng đồng thời có nhu cầu giao và nhận hàng, mỗi xe có giới hạn tải trọng và mỗi khách hàng có cửa sổ thời gian phục vụ. Hàm mục tiêu là
\[
  \text{cost} = 2000 \times \mathit{NV} + \mathit{TD},
\]
trong đó ưu tiên giảm số xe $\mathit{NV}$, sau đó giảm tổng quãng đường $\mathit{TD}$. Thuật toán gốc được trình bày trong bài báo PH-SHOWOA (PLOS ONE, 2026, doi: \href{https://doi.org/10.1371/journal.pone.0343262}{10.1371/journal.pone.0343262}).

\subsection{Bốn cải tiến chính}

\begin{table}[H]
\centering
\small
\begin{tabularx}{\textwidth}{@{}c L L@{}}
\toprule
\textbf{\#} & \textbf{Cải tiến} & \textbf{So với thuật toán gốc} \\
\midrule
1 & Full-GPU resident (CUDA, Python/Numba): toàn bộ vòng tiến hóa nằm trên GPU, CPU chỉ chuẩn bị dữ liệu và giải mã nghiệm & Chạy hoàn toàn trên CPU \\
2 & Khởi tạo SA\_RCRS\_GRASP: RCRS-GRASP $\rightarrow$ SA warmup nhẹ $\rightarrow$ Route Elimination $\rightarrow$ Deep Local Search & Khởi tạo bằng Simulated Annealing thuần \\
3 & Mô hình đa đảo (island model) với di cư vòng (ring migration) giữa các đảo & Một quần thể duy nhất, toàn bộ cá thể cùng hướng theo nghiệm tốt nhất toàn cục \\
4 & Deep Local Search cho toàn bộ quần thể, bổ sung toán tử Inter 2-opt* & Local search chỉ áp dụng cho cá thể tốt nhất \\
\bottomrule
\end{tabularx}
\caption{Các cải tiến so với thuật toán PH-SHOWOA gốc}
\end{table}

\subsection{Tình trạng hiện tại}

\begin{itemize}
  \item Hoàn thành tái lập thuật toán PH-SHOWOA gốc bằng Python (\texttt{src\_python\_baseline}), dùng làm baseline.
  \item Hoàn thành cài đặt phiên bản cải tiến full-GPU (\texttt{src\_python\_gpu\_SA\_RCRS\_GRASP}) với bốn cải tiến nêu tại Mục~1.1.
  \item Đang tiến hành thực nghiệm trên bộ benchmark 65 instance cho cả hai phiên bản; kết quả và các phân tích ablation sẽ được bổ sung khi thực nghiệm hoàn tất.
\end{itemize}

\subsection{Mã nguồn}

\begin{table}[H]
\centering
\small
\begin{tabularx}{\textwidth}{@{}>{\raggedright\arraybackslash}p{0.42\textwidth} L@{}}
\toprule
\textbf{Thư mục / tệp} & \textbf{Nội dung} \\
\midrule
\texttt{src\_python\_baseline/} & Phiên bản Python chạy trên CPU theo thuật toán gốc, dùng làm baseline \\
\texttt{src\_python\_gpu\_SA\_RCRS\_GRASP/} & Phiên bản chính chạy full-GPU, tổ chức theo Algorithm 1--10 của bài báo gốc và 4 mô-đun cải tiến \texttt{improvement1–4} \\
\texttt{src\_cpp\_rcrs\_grasp/}, \texttt{src\_cpp\_hybrid*/}, \texttt{src\_cpp\_gpu\_SA\_RCRS\_GRASP/} & Các phiên bản C++ (CPU/OpenMP, hybrid, GPU) phục vụ thử nghiệm \\
\texttt{kaggle/}, \texttt{run\_gpu\_sa\_rcrs\_grasp.py} & Script chạy benchmark trên nền tảng Kaggle (GPU) \\
\texttt{tests/} & Bộ kiểm thử tính khả thi của nghiệm và các kernel \\
\texttt{docs/} & Tài liệu mô tả cải tiến, thiết kế ablation và dữ liệu kết quả \\
\texttt{Paper/} & Bản thảo bài báo (LaTeX, template Springer Nature) \\
\bottomrule
\end{tabularx}
\caption{Cấu trúc mã nguồn của dự án}
\end{table}

\section{Các nghiên cứu được chấp nhận/công bố}

Đã có một công trình được công bố trong Lecture Notes in Computer Science (Springer, 2026):

\begin{itemize}[label={}]
  \item Lam Huynh Khang, Vo Thi Ngoc Chau, Tri Minh Pham. \textit{A Framework for Time-Series Anomaly Detection with Mamba Backbone}. Lecture Notes in Computer Science, Springer, 2026. DOI: \href{https://doi.org/10.1007/978-981-92-6510-7_14}{10.1007/978-981-92-6510-7\_14}.
\end{itemize}

Đề tài PH-SHOWOA (Mục~1) chưa được công bố; bản thảo đang được soạn thảo.

\section{Giải thưởng tại các cuộc thi}

Hiện chưa có giải thưởng tại các cuộc thi.

\section{Định hướng nghiên cứu Fall 2026 và Spring 2027}

Định hướng tiếp tục tập trung vào metaheuristic trên GPU cho bài toán định tuyến xe: học kỳ Fall 2026 hoàn thiện và nộp bài báo PH-SHOWOA-GPU; học kỳ Spring 2027 tìm hiểu các thuật toán heuristic khác và áp dụng vào PH-SHOWOA.

\subsection*{Fall 2026 -- Hoàn thiện PH-SHOWOA-GPU}

\begin{enumerate}
  \item Hoàn tất thực nghiệm trên toàn bộ 65 instance cho cả baseline và phiên bản GPU, 30 lần chạy cho mỗi instance.
  \item Thực hiện ablation cho từng cải tiến trên GPU, trên toàn bộ các instance 100 khách hàng.
  \item Đánh giá riêng hiệu quả của từng giai đoạn khởi tạo (RCRS-GRASP, SA warmup, Route Elimination, Deep Local Search) và của cơ chế ring migration.
  \item Kiểm định thống kê có hiệu chỉnh so sánh bội (Wilcoxon kết hợp Holm); đo mức tăng tốc của GPU so với CPU trên cùng ngân sách thế hệ.
  \item Hoàn thiện bản thảo và nộp bài báo (dự kiến theo định dạng tạp chí \textit{Expert Systems with Applications}).
\end{enumerate}

\subsection*{Spring 2027 -- Áp dụng các heuristic khác vào PH-SHOWOA}

\begin{enumerate}
  \item Tìm hiểu các metaheuristic khác cho bài toán VRP: Grey Wolf Optimizer (GWO), Harris Hawks Optimization (HHO), Ant Colony Optimization, Tabu Search, Variable Neighborhood Search, Iterated Local Search, ALNS và Hybrid Genetic Search (HGS).
  \item Lựa chọn các thành phần phù hợp để lai ghép vào PH-SHOWOA, chẳng hạn: cơ chế cập nhật của GWO/HHO cho bước SHO/WOA; danh sách tabu nhằm tránh lặp nghiệm trong local search; cơ chế chọn toán tử thích nghi kiểu ALNS thay cho lịch cố định (local search mỗi 25 thế hệ, di cư mỗi 20 thế hệ); toán tử lai ghép kiểu HGS.
  \item Cài đặt các thành phần mới trên khung full-GPU hiện có và thực hiện ablation từng thành phần so với PH-SHOWOA-GPU hiện tại, trên cùng tập instance, 30 lần chạy và cùng seed.
  \item Giữ lại các thành phần cải thiện cost có ý nghĩa thống kê, làm nền tảng cho bài báo thứ hai.
\end{enumerate}

\end{document}
